\section{Six-Action Trajectory Clustering Plan}
\label{sec:six-action-trajectory-clustering}

This note defines the compact action space used to cluster model-generated
trajectory turns and summarizes the validation plan for applying it to the
macro-analysis artifacts described in \texttt{analyze\_tools/MACRO\_ANALYSIS\_NOTES.md}.

\subsection{Trajectory States}

The intended state transition is:

\begin{center}
\texttt{Analyze and Plan}
$\rightarrow$
\texttt{Execute Plan and Tool Calling}
$\rightarrow$
\texttt{Verify Results}
$\rightarrow$
\texttt{Summarize and Final Answer}.
\end{center}

If verification fails, the trajectory returns to analysis and planning, then
re-enters execution and verification. The loop continues until the current
result is verified as acceptable or the trajectory ends with a limitation.

\subsection{Action Space}

Each model-generated turn receives exactly one primary action label. The action
labels are intentionally coarse, non-overlapping, and tied to trajectory
function rather than surface keywords.

\begin{table}[h]
\centering
\small
\begin{tabular}{lll}
\toprule
Action & State & Definition \\
\midrule
Problem Framing & Analyze and Plan &
Extracts the goal, givens, constraints, options, and answer format. \\
Plan Formation & Analyze and Plan &
Decomposes the task or chooses the next solving route. \\
Solution Execution & Execute Plan &
Performs internal solving: reasoning, derivation, comparison, or calculation. \\
Tool Grounding & Execute Plan and Tool Calling &
Calls a tool or uses returned tool evidence to update the solution. \\
Result Auditing & Verify Results &
Checks, diagnoses, corrects, rejects, or confirms the current result. \\
Answer Delivery & Summarize and Final Answer &
States the decisive basis and emits the final answer or terminal limitation. \\
\bottomrule
\end{tabular}
\caption{Six-action trajectory clustering space.}
\label{tab:six-action-space}
\end{table}

\subsection{Classification Unit}

The unit of classification is one assistant-generated turn. Tool-result rows are
not model actions. They are retained as observations only when a following
assistant turn explicitly uses them. User prompts, system instructions, scaffold
text, and observation-only rows should be excluded from action-rate denominators
or marked with an auxiliary flag.

\subsection{Mapping from Earlier Labels}

The earlier fine-grained labels in the macro notes map to the six actions as
follows:

\begin{table}[h]
\centering
\small
\begin{tabular}{ll}
\toprule
Earlier label & Six-action label \\
\midrule
\texttt{parse\_task} & Problem Framing \\
\texttt{decompose}, \texttt{plan\_strategy} & Plan Formation \\
\texttt{derive}, \texttt{compute}, \texttt{compare\_options} & Solution Execution \\
\texttt{consult\_tool}, \texttt{integrate\_observation} & Tool Grounding \\
\texttt{validate\_candidate}, \texttt{revise}, \texttt{detect\_error},
\texttt{confirm\_readiness} & Result Auditing \\
\texttt{commit\_answer}, \texttt{format\_answer}, \texttt{summarize\_basis},
\texttt{report\_limitation} & Answer Delivery \\
\bottomrule
\end{tabular}
\caption{Mapping from fine-grained action labels to the compact action space.}
\label{tab:fine-to-six-action-map}
\end{table}

\subsection{Tie-Break Rules}

When one turn contains multiple signals, assign the label by the most
trajectory-specific state transition it performs:

\begin{enumerate}
\item Answer Delivery
\item Tool Grounding
\item Result Auditing
\item Solution Execution
\item Plan Formation
\item Problem Framing
\end{enumerate}

Examples:

\begin{itemize}
\item A turn that says ``let me check'' and then performs arithmetic is Solution
Execution.
\item A turn that checks option-letter mapping or rejects a prior result is
Result Auditing.
\item A tool call, or a turn that explicitly uses a returned search, shell, or
code result, is Tool Grounding.
\item A terminal boxed answer, option letter, numeric value, or terminal
limitation is Answer Delivery.
\end{itemize}

\subsection{Validation Plan}

The validation question is whether the six-action space can classify all
relevant model-generated turns without systematic ambiguity.

\begin{enumerate}
\item Extract assistant-generated turns from each trajectory. Exclude user,
system, scaffold, and observation-only rows from the primary denominator.
\item Attach local context to each turn: previous action, next action, whether a
tool result is available, whether a candidate result already exists, and whether
the turn is terminal.
\item Apply the six-action priority rules in Section~\ref{sec:six-action-trajectory-clustering}.
\item Check state-transition validity. The normal path is Problem Framing to
Plan Formation to Solution Execution or Tool Grounding to Result Auditing to
Answer Delivery. Result Auditing may loop back to Plan Formation or Problem
Framing after a failed check.
\item Flag suspicious cases: Answer Delivery followed by more solving, Tool
Grounding without a tool call or used observation, Result Auditing before any
candidate or evidence exists, and repeated low-progress turns.
\item Manually audit representative correct and wrong trajectories for each
harness and benchmark covered by the macro notes.
\item Report action rates with caveats. Tool Grounding from event logs is
precise; text-derived actions are heuristic and should be reported with
low-confidence or multi-action flags when needed.
\end{enumerate}

\subsection{Expected Outliers}

The compact action space should classify the useful model-generated content, but
several outlier types should be tracked separately:

\begin{itemize}
\item Observation-only tool-result rows. These are not model actions.
\item Scaffold or instruction text accidentally included in trajectories.
\item Long DirectLLM responses that contain several logical stages inside one
raw output. These should be segmented before action clustering.
\item Low-progress loops, repeated plans, or repeated reasoning that do not
change the trajectory state.
\item HLE keyword false positives, where generic words such as ``check'' or
``therefore'' do not correspond to the actual trajectory function.
\end{itemize}

\subsection{Conclusion}

The six-action space is suitable for macro trajectory analysis because each
action corresponds to a distinct trajectory function: frame the task, form a
plan, execute internally, ground through tools, audit the result, or deliver the
answer. It is simpler and less overlapping than the earlier
\texttt{plan}/\texttt{reason}/\texttt{tool\_use}/\texttt{verify}/\texttt{recover}/\texttt{answer}
set. The main requirement is to classify assistant turns with trajectory context,
not by keywords alone.
